\documentclass[11pt]{article}
\usepackage[spanish]{babel}
\usepackage{amsmath}
\usepackage{amssymb}
\usepackage{graphicx}
\usepackage[utf8]{inputenc}
\usepackage[T1]{fontenc}
\usepackage[margin=1in]{geometry}
\usepackage{fancyhdr}
\usepackage{hyperref}
\usepackage{listings}
\usepackage{xcolor}

% --- Configuración para mostrar código C++ ---
\definecolor{codegreen}{rgb}{0,0.6,0}
\definecolor{codegray}{rgb}{0.5,0.5,0.5}
\definecolor{codepurple}{rgb}{0.58,0,0.82}
\definecolor{backcolour}{rgb}{0.95,0.95,0.92}

\lstdefinestyle{cpstyle}{
    backgroundcolor=\color{backcolour},   
    commentstyle=\color{codegreen}\itshape,
    keywordstyle=\color{blue}\bfseries,
    numberstyle=\tiny\color{codegray},
    stringstyle=\color{codepurple},
    basicstyle=\ttfamily\footnotesize,
    breakatwhitespace=false,         
    breaklines=true,                 
    captionpos=b,                    
    keepspaces=true,                 
    numbers=left,                    
    numbersep=5pt,                  
    showspaces=false,                
    showstringspaces=false,
    showtabs=false,                  
    tabsize=4
}
\lstset{style=cpstyle}

\hypersetup{
    pdftitle={Apunte Selectivo Progcomp Uchile},
    pdfauthor={Felipe Colli Olea},
    colorlinks=true,
    linkcolor=black,
    urlcolor=cyan
}

\pagestyle{fancy}
\fancyhf{}
\fancyfoot[C]{\footnotesize Felipe Colli, Todos los Derechos Reservados}
\fancyfoot[R]{\thepage}
\renewcommand{\headrulewidth}{0pt}

\title{\textbf{Apunte Selectivo Progcomp Uchile / ICPC}}
\author{Felipe Colli Olea}
\date{\today}

\begin{document}

\maketitle
\thispagestyle{fancy}
\tableofcontents
\newpage

\section{Plantillas Base}

\subsection{Template Básico (con casos de prueba)}
\begin{lstlisting}[language=C++]
#include <bits/stdc++.h>
using namespace std;

typedef long long ll;
typedef string str;
#define vec vector

int main() {
    ios_base::sync_with_stdio(0); cin.tie(0);
    int t;
    cin >> t;
    while (t--) {
        // Logica aqui
    }
    return 0;
}
\end{lstlisting}

\section{Estructuras de Datos}

\subsection{Segment Tree (Min Query)}
\begin{lstlisting}[language=C++]
struct SegmentTree {
    int n;
    vec<ll> Tree;

    SegmentTree(vec<ll>& a) {
        n = a.size();
        Tree.resize(4 * n);
        build(a, 0, 0, n - 1);
    }

    void build(vec<ll>& a, int nodo, int izq, int der) {
        if (izq == der) {
            Tree[nodo] = a[izq];
            return;
        }
        int mid = izq + (der - izq) / 2;
        build(a, 2 * nodo + 1, izq, mid);
        build(a, 2 * nodo + 2, mid + 1, der);
        Tree[nodo] = min(Tree[2 * nodo + 1], Tree[2 * nodo + 2]);
    }

    ll q(int l, int r) {
        return query(0, 0, n - 1, l, r);
    }

    ll query(int nodo, int izq, int der, int l, int r) {
        if (r < izq || l > der) return LLONG_MAX;
        if (l <= izq && der <= r) return Tree[nodo];
        int mid = izq + (der - izq) / 2;
        return min(query(2 * nodo + 1, izq, mid, l, r),
                   query(2 * nodo + 2, mid + 1, der, l, r));
    }
};
\end{lstlisting}

\subsection{Segment Tree (Max Query)}
Para Max Query, usa el mismo código de arriba, pero reemplaza \texttt{min} por \texttt{max} y el retorno fuera de rango por \texttt{LLONG\_MIN}.

\section{Grafos (Tu Salvador)}
\textit{Nota: Si los grafos no son lo tuyo, solo copia y adapta. Están optimizados y resumidos respecto al PDF para que no falles en la implementación.}

\subsection{Búsqueda en Profundidad (DFS)}
Ideal para contar componentes conexas, detectar ciclos, o simplemente recorrer.
\begin{lstlisting}[language=C++]
vector<int> visitados;
vector<vector<int>> AdjList;

void dfs(int u) {
    visitados[u] = 1;
    for (int v : AdjList[u]) {
        if (!visitados[v]) {
            dfs(v);
        }
    }
}
// En main(): visitados.assign(N, 0); AdjList.assign(N, vector<int>());
\end{lstlisting}

\subsection{Búsqueda en Anchura (BFS)}
Encuentra el \textbf{camino mínimo en grafos SIN peso} (distancia, cantidad de saltos).
\begin{lstlisting}[language=C++]
vector<int> dist;
vector<vector<int>> AdjList;

void bfs(int inicio) {
    dist[inicio] = 0;
    queue<int> q;
    q.push(inicio);
    
    while (!q.empty()) {
        int u = q.front();
        q.pop();
        
        for (int v : AdjList[u]) {
            if (dist[v] == -1) { // -1 significa no visitado
                dist[v] = dist[u] + 1;
                q.push(v);
            }
        }
    }
}
// En main(): dist.assign(N, -1);
\end{lstlisting}

\subsection{Grafos Implícitos (Grillas y Laberintos)}
Cuando el problema es un tablero/matriz, no construimos un \texttt{AdjList}. Usamos los arreglos \texttt{dr} y \texttt{dc} para movernos (Arriba, Abajo, Izq, Der).
\begin{lstlisting}[language=C++]
int MAX_FILAS, MAX_COLUMNAS;
vector<vector<char>> grilla;

// Movimientos: S, SE, E, NE, N, NW, W, SW (si piden 8 direcciones)
// Para 4 direcciones usa solo los valores correspondientes.
int dr[] = { 1, 1, 0, -1, -1, -1, 0, 1 }; 
int dc[] = { 0, 1, 1, 1, 0, -1, -1, -1 };

bool limites(int f, int c) {
    return (f >= 0 && f < MAX_FILAS && c >= 0 && c < MAX_COLUMNAS);
}

// Ejemplo explorando con BFS
void bfs_grilla(int start_f, int start_c) {
    queue<pair<int, int>> q;
    q.push({start_f, start_c});
    
    while (!q.empty()) {
        auto [f, c] = q.front(); q.pop();
        
        for (int i = 0; i < 8; ++i) { // O i < 4 si son 4 direcciones
            int nf = f + dr[i];
            int nc = c + dc[i];
            if (limites(nf, nc) /* && condicion_extra */) {
                // Procesar vecino nf, nc
                q.push({nf, nc});
            }
        }
    }
}
\end{lstlisting}

\subsection{Dijkstra (Camino Mínimo CON peso positivo)}
Encuentra la ruta más barata. ¡Ojo! Acá se usa una \texttt{priority\_queue} en reversa (de menor a mayor).
\begin{lstlisting}[language=C++]
vector<long long> dist;
// AdjList guarda {nodo_destino, peso_arista}
vector<vector<pair<int, int>>> AdjList; 
const long long INF = LLONG_MAX;

void dijkstra(int inicio) {
    dist[inicio] = 0;
    // La cola guarda {distancia_acumulada, nodo_actual}
    priority_queue<pair<long long, int>, 
                   vector<pair<long long, int>>, 
                   greater<pair<long long, int>>> pq;
                   
    pq.push({0, inicio});
    
    while (!pq.empty()) {
        auto [d_actual, u] = pq.top();
        pq.pop();
        
        // Si encontramos un camino peor guardado en la cola, lo ignoramos
        if (d_actual > dist[u]) continue;
        
        for (auto edge : AdjList[u]) {
            int v = edge.first;
            int peso = edge.second;
            
            if (dist[u] + peso < dist[v]) {
                dist[v] = dist[u] + peso;
                pq.push({dist[v], v});
            }
        }
    }
}
// En main(): dist.assign(N, INF);
\end{lstlisting}


\section{Strings y Textos}

\subsection{KMP Prefix Function}
Encuentra ocurrencias de patrones en tiempo lineal.
\begin{lstlisting}[language=C++]
vector<ll> prefix_function(const string& s) {
    int n = s.size();
    vector<ll> pi(n, 0);
    for (int i = 1; i < n; i++) {
        int j = pi[i - 1];
        while (j > 0 && s[i] != s[j])
            j = pi[j - 1];
        if (s[i] == s[j])
            j++;
        pi[i] = j;
    }
    return pi;
}
\end{lstlisting}

\subsection{Partir (Split) un String por Delimitadores}
Usando la librería de C (\texttt{<cstring>}) para separar oraciones por espacios o comas.
\begin{lstlisting}[language=C++]
string texto = "Este texto, tiene comas!";
char* token = strtok(&texto[0], " ,!"); // Delimitadores: espacio, coma, exclamacion

vector<string> palabras;
while (token != NULL) {
    palabras.push_back(token);
    token = strtok(NULL, " ,!"); 
}
\end{lstlisting}

\section{Programación Dinámica (DP)}

\subsection{Distancia de Edición (Edit Distance)}
Mínimo costo de conversiones (insertar, borrar, reemplazar) para igualar dos strings. Costo unitario = 1.
\begin{lstlisting}[language=C++]
int min3(int x, int y, int z) { return min(min(x, y), z); }

int edit_distance(string str1, string str2) {
    int m = str1.size(), n = str2.size();
    vector<vector<int>> dp(m + 1, vector<int>(n + 1));
    for (int i = 0; i <= m; i++) {
        for (int j = 0; j <= n; j++) {
            if (i == 0) dp[i][j] = j;
            else if (j == 0) dp[i][j] = i;
            else if (str1[i - 1] == str2[j - 1]) dp[i][j] = dp[i - 1][j - 1];
            else dp[i][j] = 1 + min3(dp[i][j - 1],    // Insertar
                                     dp[i - 1][j],    // Remover
                                     dp[i - 1][j - 1]); // Reemplazar
        }
    }
    return dp[m][n];
}
\end{lstlisting}

\subsection{Subsecuencia Común Más Larga (LCS)}
Encuentra el tamaño de la cadena común más larga (no necesitan ser continuas).
\begin{lstlisting}[language=C++]
int LCS(string str1, string str2) {
    int m = str1.size(), n = str2.size();
    vector<vector<int>> L(m + 1, vector<int>(n + 1, 0));
    for (int i = 0; i <= m; i++) {
        for (int j = 0; j <= n; j++) {
            if (i == 0 || j == 0) L[i][j] = 0;
            else if (str1[i - 1] == str2[j - 1])
                L[i][j] = L[i - 1][j - 1] + 1;
            else
                L[i][j] = max(L[i - 1][j], L[i][j - 1]);
        }
    }
    return L[m][n];
}
\end{lstlisting}

\section{Matemáticas / Utilidades}
\subsection{Generar todas las permutaciones}
\begin{lstlisting}[language=C++]
string s = "CBA";
sort(s.begin(), s.end()); // IMPORTANTISIMO: Siempre ordenar antes
vector<string> salida;
do {
    salida.push_back(s);
} while (next_permutation(s.begin(), s.end()));
\end{lstlisting}


\section{Estructuras "Mágicas" y Fáciles}

\subsection{Ordered Set (PBDS)}
Es un \texttt{set} normal, pero tiene dos funciones extra: \texttt{find\_by\_order(k)} (devuelve iterador al k-ésimo elemento, 0-indexado) y \texttt{order\_of\_key(x)} (cuenta cuántos elementos son estrictamente menores a x).
\begin{lstlisting}[language=C++]
#include <ext/pb_ds/assoc_container.hpp>
#include <ext/pb_ds/tree_policy.hpp>
using namespace __gnu_pbds;

template <typename T>
using ordered_set = tree<T, null_type, less<T>, rb_tree_tag, tree_order_statistics_node_update>;

// Uso:
// ordered_set<int> st;
// st.insert(5); st.insert(2);
// cout << *st.find_by_order(0); // Imprime 2
// cout << st.order_of_key(5);   // Imprime 1 (solo el 2 es menor)
\end{lstlisting}

\subsection{Fenwick Tree (BIT)}
Suma en rangos y actualización de un punto en $O(\log n)$. Más corto que SegTree. (1-indexado internamente).
\begin{lstlisting}[language=C++]
struct BIT {
    int n; vector<ll> bit;
    BIT(int n) : n(n), bit(n + 1, 0) {}
    void add(int idx, ll val) {
        for (++idx; idx <= n; idx += idx & -idx) bit[idx] += val;
    }
    ll query(int idx) { // suma de [0, idx]
        ll ret = 0;
        for (++idx; idx > 0; idx -= idx & -idx) ret += bit[idx];
        return ret;
    }
    ll query(int l, int r) { return query(r) - query(l - 1); }
};
\end{lstlisting}

\subsection{Disjoint Set Union (DSU)}
Agrupa elementos y verifica si dos elementos están en el mismo grupo. Excelente para problemas de conectividad o Kruskal (MST).
\begin{lstlisting}[language=C++]
struct DSU {
    vector<int> p, sz;
    DSU(int n) {
        p.resize(n); iota(p.begin(), p.end(), 0);
        sz.assign(n, 1);
    }
    int find(int i) {
        if (p[i] == i) return i;
        return p[i] = find(p[i]); // Path compression
    }
    bool unite(int i, int j) {
        int root_i = find(i), root_j = find(j);
        if (root_i == root_j) return false;
        if (sz[root_i] < sz[root_j]) swap(root_i, root_j);
        p[root_j] = root_i;
        sz[root_i] += sz[root_j];
        return true;
    }
};
\end{lstlisting}


\section{Matemáticas}

\subsection{Criba de Eratóstenes (Sieve)}
Encuentra todos los primos hasta $N$ y el factor primo más pequeño (SPF) de cada número.
\begin{lstlisting}[language=C++]
const int MAXN = 1e6 + 5; // Cambiar segun problema
int spf[MAXN]; // Smallest Prime Factor
vector<int> primes;

void sieve() {
    for (int i = 2; i < MAXN; i++) spf[i] = i;
    for (int i = 2; i * i < MAXN; i++) {
        if (spf[i] == i) {
            for (int j = i * i; j < MAXN; j += i)
                if (spf[j] == j) spf[j] = i;
        }
    }
    for (int i = 2; i < MAXN; i++) 
        if (spf[i] == i) primes.push_back(i);
}
// factorizar rapido: while(n > 1) { factors.push_back(spf[n]); n /= spf[n]; }
\end{lstlisting}

\subsection{Exponenciación Binaria y Modulo}
Calcula $a^b \pmod m$ en $O(\log b)$.
\begin{lstlisting}[language=C++]
ll binpow(ll a, ll b, ll m) {
    a %= m;
    ll res = 1;
    while (b > 0) {
        if (b & 1) res = res * a % m;
        a = a * a % m;
        b >>= 1;
    }
    return res;
}

// Inverso modular de A mod M (Solo funciona si M es primo)
ll modInverse(ll a, ll m) {
    return binpow(a, m - 2, m);
}
// Division modular (a / b) mod m
ll modDivide(ll a, ll b, ll m) {
    return (a % m) * modInverse(b, m) % m;
}
\end{lstlisting}

\end{document}
